\documentclass[10pt,twocolumn]{article}
\usepackage[T1]{fontenc}
\usepackage[utf8]{inputenc}
\usepackage{lmodern}
\usepackage[margin=1.8cm,top=2.4cm,bottom=2cm,columnsep=0.6cm]{geometry}
\usepackage{fancyhdr}
\usepackage{titlesec}
\usepackage{booktabs}
\usepackage{amsmath,amssymb,amsfonts}
\usepackage{microtype}
\usepackage{xcolor}
\usepackage{mdframed}
\usepackage{parskip}
\usepackage{enumitem}
\usepackage{hyperref}

\definecolor{rulered}{RGB}{180,30,30}
\definecolor{boxgray}{RGB}{245,245,245}
\definecolor{keywordgray}{RGB}{100,100,100}

\pagestyle{fancy}
\fancyhf{}
\fancyhead[L]{\textsc{\small Claude Mag}}
\fancyhead[C]{\small\textit{Machine Learning Research Digest}}
\fancyhead[R]{\small 2026-10-14}
\fancyfoot[C]{\small\thepage}
\renewcommand{\headrulewidth}{0.8pt}

\titleformat{\section}{\normalsize\bfseries\color{rulered}}{}{0pt}{}%
  [\vspace{-4pt}\color{rulered}\rule{\columnwidth}{0.4pt}]
\titlespacing{\section}{0pt}{8pt}{4pt}

\hypersetup{colorlinks=true,urlcolor=rulered,linkcolor=black}

\begin{document}
\setlength{\parindent}{0pt}
\setlength{\parskip}{4pt}

{\small\color{keywordgray}\textbf{Keywords:}
  chain-of-thought \textbullet{}
  test-time compute \textbullet{}
  expansion utility \textbullet{}
  self-consistency}\par\vspace{4pt}

{\LARGE\bfseries\raggedright
  Expanding LLM Reasoning}\par\vspace{4pt}

{\large\itshape\raggedright
  Branching From the Last Step of an Existing Chain Beats Uniform Placement
  and Rivals Learned Routers While Using 0.77$\times$ the Tokens of
  Self-Consistency}\par\vspace{6pt}

{\small\textbf{By} Rian Atri \& Evan Luo}\par\vspace{2pt}

{\small\color{keywordgray}arXiv:2610.05584 \textbullet{} NeurIPS 2026 \textbullet{} October 2026}
\par\vspace{2pt}\noindent{\color{rulered}\rule{\columnwidth}{0.6pt}}\vspace{6pt}

Instead of sampling multiple full chains from scratch, \emph{Expanding LLM
Reasoning} asks: given one completed chain-of-thought trace, at which step
should an additional continuation branch? A simple ``always branch from the
last step'' rule surpasses uniform placement across 41 model-and-benchmark
cells and uses only 0.774$\times$ the generated tokens of four-sample
self-consistency at matched accuracy on MATH-500.

\begin{mdframed}[backgroundcolor=boxgray,linecolor=rulered,linewidth=1pt,
    innerleftmargin=6pt,innerrightmargin=6pt,
    innertopmargin=6pt,innerbottommargin=6pt]
\textbf{\small AT A GLANCE}\par\vspace{4pt}
\begin{description}[leftmargin=0pt,labelsep=4pt,itemsep=2pt,parsep=0pt,
    font=\small\bfseries]
  \item[Problem:] Self-consistency generates full chains from scratch; it
    ignores information already encoded in a completed trace.
  \item[Why it matters:] Inference compute is the bottleneck for large
    reasoning models; reusing partial traces can reduce cost.
  \item[Method:] Define expansion utility at each step; compare uniform,
    always-last, and learned-router branching policies across 41 cells.
  \item[Result:] Always-last matches self-consistency accuracy using
    0.774$\times$ the tokens on MATH-500; $+$4.25-point oracle headroom
    motivates future learned selectors.
  \item[Evidence:] Nine models, six benchmarks, cross-fitted held-out
    evaluation.
\end{description}
\end{mdframed}

\section{The Big Picture}

Chain-of-thought (CoT) reasoning substantially improves LLM accuracy on hard
tasks, but generating complete independent chains is expensive. Self-consistency
and best-of-$n$ approaches sample $n$ full chains, consuming $n\times$ the
output tokens.

This paper reframes the problem: if one chain $\tau$ is already in hand, at which
step $t^*$ should a continuation branch to maximally improve the probability of
reaching the correct answer? The key insight is that later steps encode more
context and have lower continuation variance---making the last step a reliable
default branching point regardless of the task or model.

\section{Why Existing Approaches Fall Short}

\textbf{Self-consistency} samples complete chains from scratch and majority-votes.
It discards the partial trace information already computed.

\textbf{Best-of-$n$ with process reward models} requires training a step-level
critic, which is expensive and brittle across domains.

\textbf{Learned branching routers} require training data and do not consistently
beat a simple positional heuristic, as this paper shows.

\section{Core Mechanism: Expansion Utility}

The authors define \textbf{expansion utility} at step $t$ as the expected
correctness improvement from branching at $t$:
\[
U(t) = \mathbb{E}_{c \sim p(\cdot|\tau_{:t})}[\mathbf{1}[a(c)=a^*]]
       - \mathbf{1}[a(\tau)=a^*]
\]
Positive $U(t)$ means branching at $t$ is expected to recover a failed chain.

Three branching policies are compared:
\begin{itemize}[itemsep=2pt,parsep=0pt]
  \item \textbf{Uniform:} branch at a randomly selected step.
  \item \textbf{Always-last:} always branch at the final step $|\tau|$.
  \item \textbf{Learned router:} a 2-layer MLP predicting $U(t)>0$ from
    prefix hidden states.
\end{itemize}

\section{Technical Formulation}

Let $\tau = (s_1, \ldots, s_T)$ be a chain and $\tau_{:t}$ its prefix.
With $k$ continuations sampled from step $t$, the expected accuracy is:
\[
\mathrm{Acc}(t,k) = \mathbb{E}_{c_1,\ldots,c_k \sim p(\cdot|\tau_{:t})}
\!\left[\mathrm{MajVote}(a(c_1),\ldots,a(c_k))=a^*\right]
\]

The optimal policy selects $t^* = \arg\max_t \mathrm{Acc}(t,k)$.

For the learned router, let $\phi(\tau_{:t}) \in \mathbb{R}^d$ be the model's
last-layer hidden state at step $t$. The router solves:
\[
\hat{y}(t) = \mathbf{1}[\mathbf{w}^\top \phi(\tau_{:t}) + b > 0]
\]
trained with binary cross-entropy on labels $\mathbf{1}[U(t)>0]$.

Held-out expansion utility at selected step $t^*$:
\[
\Delta\mathrm{Acc}(\pi) = \mathrm{Acc}(\pi(t), k) - \mathrm{Acc}(\mathrm{Uniform}, k)
\]
measures the gain of policy $\pi$ over the uniform baseline.

\section{Learning or Inference Procedure}

\begin{enumerate}[itemsep=2pt,parsep=0pt]
  \item Generate one chain $\tau$ per problem instance.
  \item Apply the branching policy to select step $t^*$.
  \item Sample $k{=}4$ continuations from $\tau_{:t^*}$ in parallel.
  \item Majority-vote over $k$ answers.
  \item Compare total tokens against self-consistency at matched accuracy.
\end{enumerate}

The router is trained once per model on 1,000 labeled chains; evaluation
uses cross-fitting (router trained on split A, evaluated on split B).

\section{What the Guarantee Says}

The paper establishes empirical monotonicity of variance across steps: later
steps show lower continuation-correctness variance than earlier steps across
38/41 model-benchmark cells. This explains why always-last is a robust heuristic
even without a learned predictor.

Oracle selection outperforms always-last by $+$4.25 points on average, showing
residual headroom that better future selectors could capture.

\section{Experimental Findings}

\begin{itemize}[itemsep=2pt,parsep=0pt]
  \item \textbf{Models:} DeepSeek-R1-Distill-Qwen-14B, Qwen3-14B, Gemma-2-9B,
    QwQ-32B, five smaller models.
  \item \textbf{Benchmarks:} MATH-500, AIME 2024/25, AMC 2023, GPQA-Diamond,
    LiveCodeBench.
  \item Always-last beats self-consistency at \textbf{0.774$\times$ tokens}
    on MATH-500 and 0.81$\times$ on AIME 2024.
  \item Learned router shows $+$4.25 gain on held-out cross-fitted sets vs.\
    uniform, but not over always-last.
  \item Always-last beats uniform in 38/41 cells; ties in 3.
\end{itemize}

\section{Ablations and Interpretation}

\textbf{Branch depth:} Last vs.\ second-to-last step gives similar accuracy;
the last step carries the most cumulative context.

\textbf{Continuation length:} Full-chain continuations outperform few-step
endings slightly, suggesting the model benefits from regenerating the answer tokens.

\textbf{Number of branches $k$:} Accuracy improves log-linearly up to $k{=}8$;
diminishing returns beyond.

\textbf{Router architecture:} 2-layer MLP outperforms linear probes; 4+ layers
overfit the training split.

\textbf{Why always-last works:} Empirical $U(t)$ variance decreases monotonically
with step depth; the last step minimises continuation variance, making it the
most stable branching point.

\section*{Reference}

Rian Atri, Evan Luo.
\textbf{Expanding LLM Reasoning}.
NeurIPS 2026. arXiv:2610.05584, October 2026.\\
\url{https://arxiv.org/abs/2610.05584}

\end{document}
